% Options for packages loaded elsewhere
\PassOptionsToPackage{unicode}{hyperref}
\PassOptionsToPackage{hyphens}{url}
%
\documentclass[
  10pt,
]{article}
\usepackage{amsmath,amssymb}
\usepackage{iftex}
\ifPDFTeX
  \usepackage[T1]{fontenc}
  \usepackage[utf8]{inputenc}
  \usepackage{textcomp} % provide euro and other symbols
\else % if luatex or xetex
  \usepackage{unicode-math} % this also loads fontspec
  \defaultfontfeatures{Scale=MatchLowercase}
  \defaultfontfeatures[\rmfamily]{Ligatures=TeX,Scale=1}
\fi
\usepackage{lmodern}
\ifPDFTeX\else
  % xetex/luatex font selection
\fi
% Use upquote if available, for straight quotes in verbatim environments
\IfFileExists{upquote.sty}{\usepackage{upquote}}{}
\IfFileExists{microtype.sty}{% use microtype if available
  \usepackage[]{microtype}
  \UseMicrotypeSet[protrusion]{basicmath} % disable protrusion for tt fonts
}{}
\makeatletter
\@ifundefined{KOMAClassName}{% if non-KOMA class
  \IfFileExists{parskip.sty}{%
    \usepackage{parskip}
  }{% else
    \setlength{\parindent}{0pt}
    \setlength{\parskip}{6pt plus 2pt minus 1pt}}
}{% if KOMA class
  \KOMAoptions{parskip=half}}
\makeatother
\usepackage{xcolor}
\usepackage[margin=0.85in]{geometry}
\usepackage{longtable,booktabs,array}
\usepackage{calc} % for calculating minipage widths
% Correct order of tables after \paragraph or \subparagraph
\usepackage{etoolbox}
\makeatletter
\patchcmd\longtable{\par}{\if@noskipsec\mbox{}\fi\par}{}{}
\makeatother
% Allow footnotes in longtable head/foot
\IfFileExists{footnotehyper.sty}{\usepackage{footnotehyper}}{\usepackage{footnote}}
\makesavenoteenv{longtable}
\setlength{\emergencystretch}{3em} % prevent overfull lines
\providecommand{\tightlist}{%
  \setlength{\itemsep}{0pt}\setlength{\parskip}{0pt}}
\setcounter{secnumdepth}{-\maxdimen} % remove section numbering
\ifLuaTeX
  \usepackage{selnolig}  % disable illegal ligatures
\fi
\IfFileExists{bookmark.sty}{\usepackage{bookmark}}{\usepackage{hyperref}}
\IfFileExists{xurl.sty}{\usepackage{xurl}}{} % add URL line breaks if available
\urlstyle{same}
\hypersetup{
  pdftitle={Recoverable Feature Routing under Distribution Shift: Fragile dominance and few-shot recovery},
  pdfauthor={Mykyta Siryk},
  hidelinks,
  pdfcreator={LaTeX via pandoc}}

\title{Recoverable Feature Routing under Distribution Shift: Fragile
dominance and few-shot recovery}
\author{Mykyta Siryk}
\date{4 October 2026, 18:46 UTC}

\begin{document}
\maketitle

\hypertarget{abstract}{%
\subsubsection{Abstract}\label{abstract}}

A dominant but fragile eye cue is separated from weaker lower-face and
chromatic evidence. Held-out glasses reduce fixed global fusion to
0.5000 accuracy. Across five seeds, few-shot routing reaches 0.6302 +/-
0.1064 at eight shots and shifts routing mass toward backup evidence,
but some seeds remain locked. The experiment demonstrates retention and
partial recoverability while exposing high variance.

\hypertarget{historical-manuscript-status}{%
\subsubsection{Historical Manuscript
Status}\label{historical-manuscript-status}}

This is the cumulative scientific manuscript at release \textbf{v0.2.0}.
It includes only experiments completed by 4 October 2026, 18:46 UTC;
later findings are intentionally absent.

\hypertarget{introduction}{%
\subsubsection{1. Introduction}\label{introduction}}

Neural classifiers usually prioritize features that are most predictive
across their training distribution. That behavior is efficient when
feature reliability is stable, but brittle when a dominant cue fails
while lower-priority evidence remains informative. The motivating
deepfake-detection case is a model that trusts one family of texture,
frequency, chromatic, geometric, or localized facial evidence until
eyewear, occlusion, compression, or generator shift changes its
reliability.

The desired response is not necessarily permanent retraining. A system
should reconsider which evidence it trusts for the affected regime while
preserving its established behavior elsewhere. This manuscript records a
sequence of controlled falsification experiments designed to separate
apparent adaptation from actual access to neglected evidence.

Each tagged version is a historical scientific snapshot. Negative
results, implementation failures, and abandoned explanations are
retained because they determine why the next experiment exists.

\hypertarget{related-work}{%
\subsubsection{2. Related Work}\label{related-work}}

Mixture-of-experts systems use learned gates for input-dependent expert
selection {[}1,2{]}. MAML and fast-weight systems formalize rapid
adaptation {[}3,4{]}, while test-time adaptation and continual learning
expose complementary distribution-shift and stability problems
{[}5,6{]}. The present sequence isolates a narrower question: after a
learned evidence policy becomes unreliable, does useful backup evidence
remain both represented and reachable by a bounded adaptation process?

\hypertarget{controlled-experimental-setting}{%
\subsubsection{3. Controlled Experimental
Setting}\label{controlled-experimental-setting}}

\hypertarget{evidence-families}{%
\paragraph{3.1 Evidence families}\label{evidence-families}}

Each synthetic image belongs to one of two classes and contains three
separately accessible evidence families:

\begin{itemize}
\tightlist
\item
  \textbf{A: eye-region evidence.} A nearly deterministic localized
  marker makes this the dominant cue on clean data.
\item
  \textbf{B: lower-face geometry.} A weaker geometric pattern is
  correlated with the label with probability 0.72.
\item
  \textbf{C: chromatic evidence.} A weaker global color bias is
  correlated with the label with probability 0.66.
\end{itemize}

Three small convolutional experts receive feature-isolated views. Expert
A sees the eye band, expert B sees grayscale lower-face geometry, and
expert C sees an opponent-color representation.

The separation is intentional. It allows us to measure whether a model
retains genuinely useful backup information after A fails, rather than
inferring feature behavior from a monolithic latent representation.

\hypertarget{held-out-intervention}{%
\paragraph{3.2 Held-out intervention}\label{held-out-intervention}}

The final intervention is a synthetic glasses overlay applied to the eye
band. It is not used during slow-router training. The perturbation is
designed to invalidate A while leaving B and C largely intact.

The router is meta-trained on other perturbations, including eye-band
masking, lower-face masking, desaturation, blur, JPEG compression, and
clean inputs. Thus, the system is exposed to the general idea that
evidence reliability can change, but it does not learn a direct
glasses-specific rule.

\hypertarget{evaluation-protocol}{%
\paragraph{3.3 Evaluation protocol}\label{evaluation-protocol}}

The principal runs use five seeds: 11, 22, 33, 44, and 55. Few-shot
support sizes are 2, 4, 8, 16, and 32 labeled glasses examples. Support
samples are excluded from query evaluation.

The main metrics are shifted-regime accuracy, clean accuracy after
adaptation, expert routing weights, routing entropy, and saturation
rate. Later experiments add explicit measures of dominant-cue release,
backup-cue recovery, routing Jacobian conductance, and initial
backup-gradient magnitude.

\hypertarget{experiments}{%
\subsubsection{4. Experiments}\label{experiments}}

\hypertarget{experiment-1-failure-of-an-overly-easy-toy-world}{%
\paragraph{4.1 Experiment 1: failure of an overly easy toy
world}\label{experiment-1-failure-of-an-overly-easy-toy-world}}

The first toy construction contained RGB, edge, high-pass, and chromatic
expert views. It failed as a falsification environment because all
fusion strategies achieved 1.000 accuracy on both clean and
glasses-corrupted data. The perturbation increased cross-entropy loss,
but the edge expert remained sufficiently diagnostic that the router had
no incentive to rediscover weak evidence.

This negative result motivated a stricter construction in which the
globally dominant feature is explicitly fragile.

\hypertarget{experiment-2-fragile-dominance-and-few-shot-recovery}{%
\paragraph{4.2 Experiment 2: fragile dominance and few-shot
recovery}\label{experiment-2-fragile-dominance-and-few-shot-recovery}}

Experiment 2 introduced evidence families A, B, and C. On clean data,
the learned global fusion placed approximately 95.7\% of routing mass on
A. Under glasses, the fixed global system fell to exactly 0.5000 mean
accuracy across five seeds.

The backup experts remained useful. Under glasses, B and C independently
retained above-chance classification performance, while A collapsed to
chance. A support-optimized fusion reference reached 0.6944 ± 0.0302
accuracy at 32 shots.

The adaptive router partially recovered:

\begin{longtable}[]{@{}lrr@{}}
\toprule\noalign{}
Method & K & Glasses accuracy \\
\midrule\noalign{}
\endhead
\bottomrule\noalign{}
\endlastfoot
Global fixed fusion & 0 & 0.5000 ± 0.0000 \\
Adaptive zero-shot & 0 & 0.5008 ± 0.0016 \\
Adaptive few-shot & 2 & 0.5872 ± 0.1052 \\
Adaptive few-shot & 8 & \textbf{0.6302 ± 0.1064} \\
Adaptive few-shot & 32 & 0.6300 ± 0.1061 \\
Support-optimized fusion reference & 32 & 0.6944 ± 0.0302 \\
\end{longtable}

At 8 shots, mean routing changed from approximately

\[
[A,B,C]=[0.904,0.088,0.008]
\]

before glasses-specific adaptation to

\[
[A,B,C]=[0.451,0.539,0.010].
\]

The mechanism therefore worked in aggregate, but the standard deviation
was large. Inspection showed that some seeds escaped from A while others
remained almost perfectly locked to it.

\hypertarget{interpretation-at-this-milestone}{%
\subsubsection{Interpretation at This
Milestone}\label{interpretation-at-this-milestone}}

Backup evidence remains represented and is sometimes recoverable, but
access is seed-dependent. Feature presence and reliable rapid adaptation
are distinct properties.

\hypertarget{reason-for-the-next-stage}{%
\subsubsection{Reason for the Next
Stage}\label{reason-for-the-next-stage}}

The limitations and unresolved mechanism at this milestone directly
motivate Experiment 3; the next tagged manuscript records that test
rather than silently rewriting this version.

\hypertarget{limitations-and-threats-to-validity}{%
\subsubsection{Limitations and Threats to
Validity}\label{limitations-and-threats-to-validity}}

The controlled worlds are binary, visually regular, and explicitly
factorized. The support sets are labeled, only five principal seeds are
used after the first stages, and the synthetic glasses intervention may
be easier to recognize than natural distribution shifts.
Best-operating-point comparisons may benefit from multiple testing.
Experiment 0 lacks raw result artifacts, while Experiment 6 lacks
completed results. Executable smoke tests establish software integrity,
not reproduction of full-scale scientific outcomes.

\hypertarget{reproducibility-and-provenance}{%
\subsubsection{Reproducibility and
Provenance}\label{reproducibility-and-provenance}}

The accompanying repository separates byte-preserved historical source
and results from repaired runnable snapshots. Each runnable stage
provides a \texttt{uv} environment, exact entrypoint, bounded smoke
test, and validation record. Reconstructed proposal, implementation,
results, and analysis commits preserve why the sequence changed.
Numerical claims come from stored outputs unless explicitly marked
reported-only or results-missing.

\hypertarget{disclosure-of-ai-assistance}{%
\subsubsection{Disclosure of AI
Assistance}\label{disclosure-of-ai-assistance}}

ChatGPT collaborated in experimental ideation, implementation drafting,
debugging, interpretation, manuscript preparation, and archival
reconstruction. Git commits credit ChatGPT through co-author trailers.
The exact ChatGPT model identifiers and reasoning-effort settings used
during the original 4 October 2026 experiment conversations were not
exposed in the recovered metadata and have not been inferred. Mykyta
Siryk remains responsible for reviewing the code, evidence, claims, and
manuscript.

\hypertarget{references}{%
\subsubsection{References}\label{references}}

{[}1{]} N. Shazeer, A. Mirhoseini, K. Maziarz, A. Davis, Q. V. Le, G. E.
Hinton, and J. Dean, ``Outrageously Large Neural Networks: The
Sparsely-Gated Mixture-of-Experts Layer,'' ICLR, 2017.

{[}2{]} W. Fedus, B. Zoph, and N. Shazeer, ``Switch Transformers:
Scaling to Trillion Parameter Models with Simple and Efficient
Sparsity,'' Journal of Machine Learning Research, vol.~23, no. 120,
pp.~1--39, 2022.

{[}3{]} C. Finn, P. Abbeel, and S. Levine, ``Model-Agnostic
Meta-Learning for Fast Adaptation of Deep Networks,'' Proceedings of the
34th International Conference on Machine Learning, PMLR 70,
pp.~1126--1135, 2017.

{[}4{]} J. Ba, G. Hinton, V. Mnih, J. Z. Leibo, and C. Ionescu, ``Using
Fast Weights to Attend to the Recent Past,'' Advances in Neural
Information Processing Systems, 2016.

{[}5{]} D. Wang, E. Shelhamer, S. Liu, B. Olshausen, and T. Darrell,
``Tent: Fully Test-Time Adaptation by Entropy Minimization,''
International Conference on Learning Representations, 2021.

{[}6{]} J. Kirkpatrick et al., ``Overcoming Catastrophic Forgetting in
Neural Networks,'' Proceedings of the National Academy of Sciences,
vol.~114, no. 13, pp.~3521--3526, 2017.

\end{document}
